\clearpage

\begin{figure}[p]
    \caption{Gold price and exchange rates from 1931 to 1936}
    \label{fig:F2}
    
    \begin{center} Panel A \end{center}
    \makebox[\textwidth][c]{\includegraphics[width=0.56\textwidth,height = 6.79cm]{figures/body/dollar_to_sterling.pdf}}
    
    \vspace*{0.5cm}
    \begin{center} Panel B \end{center}
    \makebox[\textwidth][c]{\includegraphics[width=0.56\textwidth,height = 6.79cm]{figures/body/dollar_to_franc.pdf}}
    \vspace*{0.75cm}
    
    \footnotesize
    \emph{Notes}: The red-dashed line represents the dollar/sterling (Panel A) and dollar/franc (Panel B) exchange rates normalized to 1 in December 1932. The black vertical line represents the start of the requirement to return all gold holdings to the U.S. government. The solid-blue line represents the gold price in the U.S. government's purchasing programs. All variables are at a monthly frequency.  
\end{figure}

\clearpage

\begin{figure}[p]
    \caption{Gold price and inflation in 1933 and 1934}
	\label{fig:F3}
	
	\makebox[\textwidth][c]{\includegraphics[width=0.56\textwidth,height = 6.79cm]{figures/body/inflation.pdf}}
    
	\vspace*{1cm}
	\footnotesize
	\emph{Notes}: The blue-dotted line represents the official domestic gold price in the United States. The vertical black line represents the start of the requirement to return all gold holdings to the government in the United States. The solid-blue line is the gold price in the  U.S. government's purchasing programs. The red-dashed line is the Consumer Price Index (CPI) normalized to 1 in December 1932.
\end{figure}

\clearpage

\begin{figure}[p]
    \caption{Gold clause effect on investment by year}
    \begin{center}
    \includegraphics[width=0.75\textwidth, height = 8.79cm]{figures/body/parallel_trend_plot.pdf}
    \end{center}
    \footnotesize \textit{Notes.} This figure plots the year dummy $\times$ $\tilde{d}$ interaction coefficients and 95\% confidence bands from the investment regression reported in column 2 of Table \ref{tab:inv_main}. The specification regresses net investment on $Q$ and year dummy $\times$ $\tilde{d}$ interactions with firm and year fixed effects.
    \label{fig:gold_coeffs}
\end{figure}

\clearpage

\begin{figure}[p]
    \caption{New bond offerings, industrial firms (Hickman, 1953)}
    \begin{center}
    \includegraphics[width=0.75\textwidth, height = 8.79cm]{figures/body/industrial_corp_bond_issuance.pdf}
    \end{center}
    \footnotesize \textit{Notes.} This figure plots the issuance of industrial corporate bonds by year.
    \label{fig:bond_issuance}
\end{figure}

